% !TeX root = ../../main.tex
\subsection{T-based inference}

At convergence, the model-based covariance estimator is
\begin{equation}
  \widehat{\operatorname{Cov}}
  (\widehat{\boldsymbol{\beta}})
  =(X^\top\widehat{W}X)^{-1}.
  \label{eq:covariance}
\end{equation}
The Pearson equation scales the residual variation to approximately one.  If
the dispersion estimate reaches its numerical upper boundary, the
implementation additionally multiplies Equation~\eqref{eq:covariance} by the
remaining Pearson scale.

Equation~\eqref{eq:covariance} is a plug-in, model-based covariance: it
conditions on the estimated $\widehat{\rho}$ as though it were fixed.  It does
not propagate uncertainty from estimating a separate dispersion for every
guide.  Referring the coefficient to $t_{m-q}$ gives a heavier-tailed
small-sample reference than a normal approximation, but it is not a formal
correction for dispersion-estimation uncertainty.  With few independent
libraries, calibration must therefore be checked by simulation, phenotype
permutation, or prespecified negative controls.

For coefficient $\beta_j$, the statistic
\begin{equation}
  t_j =
  \frac{\widehat{\beta}_j-\beta_{j,0}}
  {\sqrt{
  \widehat{\operatorname{Cov}}
  (\widehat{\boldsymbol{\beta}})_{jj}}}
  \label{eq:tstat}
\end{equation}
is compared with a Student $t_{m-q}$ distribution.  The degrees of freedom are
driven by the number of independent sequenced samples, not by the total read
count.  This is the crucial reason that a library with millions of reads does
not create artificial certainty.

More generally, for a prespecified contrast vector $\mathbf{c}$,
\begin{equation}
  t_{\mathbf{c}} =
  \frac{\mathbf{c}^\top\widehat{\boldsymbol{\beta}}-\delta_0}
  {\sqrt{\mathbf{c}^\top
  \widehat{\operatorname{Cov}}(\widehat{\boldsymbol{\beta}})
  \mathbf{c}}}
\end{equation}
uses the same $m-q$ degrees of freedom.  This covers pairwise group contrasts,
average slopes in interaction models, and other one-degree-of-freedom
hypotheses.  Guide-wise two-sided $p$-values are adjusted using the
Benjamini--Hochberg procedure \citep{benjamini1995controlling}.

\subsubsection{From guide statistics to gene statistics}

BARCS fits the count model guide by guide.  The primary HT-29, FACS,
simulation, GSE70038, and A375 benchmark scripts then use an explicit
directional Stouffer summary, denoted BARCS-ST below, to obtain one gene-level
statistic.
For the $j$th valid guide targeting gene $g$, the two-sided guide $p$-value is
converted back to a signed standard-normal score,
\begin{equation}
  z_{gj}
  =
  \operatorname{sign}(\widehat{\beta}_{gj})
  \Phi^{-1}\!\left(1-\frac{p_{gj}}{2}\right).
  \label{eq:signedz}
\end{equation}
If $m_g$ valid guides target the gene, their combined score and two-sided
gene-level $p$-value are
\begin{equation}
  Z_g=\frac{1}{\sqrt{m_g}}\sum_{j=1}^{m_g}z_{gj},
  \qquad
  p_g=2\Phi(-|Z_g|).
  \label{eq:genestouffer}
\end{equation}
The reported gene effect is the median of the guide coefficient estimates,
\begin{equation}
  \widehat{\beta}_g
  =\operatorname{median}_{j=1,\ldots,m_g}
    \widehat{\beta}_{gj},
  \label{eq:geneeffect}
\end{equation}
and gene-level false-discovery rates are obtained by applying
Benjamini--Hochberg to the collection of $p_g$ values.  Consequently,
concordant guide directions reinforce one another, discordant directions
cancel, and one extreme guide has limited influence on the reported effect.

This aggregation is deliberately transparent, but it is not a native
hierarchical gene model.  Equations~\eqref{eq:signedz}--\eqref{eq:genestouffer}
give equal weight to valid guides and use the independence reference
$\sqrt{m_g}$; shared sequence artifacts or other correlation can therefore
make the gene statistic anti-conservative.  A production extension should
estimate guide reliability and within-gene dependence or use a hierarchical
effect model.  In the biological head-to-head benchmarks, this
Stouffer--median procedure is used only for BARCS (and for the deliberately
naive binomial comparator).  MAGeCK, Chronos, BAGEL2, Waterbear, and MAUDE
retain their native or deposited gene-level summaries.  The exploratory
DeWeirdt false-discovery audit used Fisher aggregation as stated in its
Results description and is not pooled with these primary benchmark
statistics.

\paragraph{Exchangeable normal guide-beta statistic.}
BARCS-NORM implements a different and deliberately simple model.  It does not
transform guide $p$-values and does not use guide-specific standard errors.
Instead, for the $m_g$ valid guides targeting gene $g$, it assumes
\begin{equation}
  \widehat{\beta}_{gj}
  \mathrel{\overset{\mathrm{iid}}{\sim}}
  N(\mu_g,\sigma_g^2).
  \label{eq:genenormal}
\end{equation}
The gene effect and between-guide standard deviation are estimated directly:
\begin{align}
  \widehat{\mu}_g
  &=
  \frac{1}{m_g}\sum_{j=1}^{m_g}\widehat{\beta}_{gj},\\
  s_g^2
  &=
  \frac{1}{m_g-1}
  \sum_{j=1}^{m_g}
  (\widehat{\beta}_{gj}-\widehat{\mu}_g)^2.
\end{align}
Testing $H_0:\mu_g=0$ gives
\begin{equation}
  T_g
  =
  \frac{\widehat{\mu}_g}{s_g/\sqrt{m_g}}
  \sim t_{m_g-1}
  \quad\text{under Equation~\eqref{eq:genenormal} and }H_0.
  \label{eq:genenormalt}
\end{equation}
Thus normally distributed guide betas lead to a Student $t$ reference when
their unknown $\sigma_g$ is estimated from the same guides.  The software also
reports $2\Phi(-|T_g|)$ as a plug-in normal sensitivity value, but it is not
the primary $p$-value because it treats the estimated spread as known.
Multiple guides support this guide-consistency calculation; they do not become
additional biological screen replicates.

\paragraph{Random-effects and moderated random-effects statistics.}
BARCS-RE (the implementation label \texttt{BARCS-partial}) retains each
guide-specific regression standard error and models
\begin{equation}
  \widehat{\beta}_{gj}\mid\beta_g,\tau_g^2
  \sim
  N\!\left(
    \beta_g,\,
    \operatorname{SE}(\widehat{\beta}_{gj})^2+\tau_g^2
  \right).
  \label{eq:generandomeffects}
\end{equation}
It estimates $\tau_g^2$ by the nonnegative DerSimonian--Laird estimator
\citep{dersimonian1986meta} and then pools guide effects with weights
\begin{equation}
  w_{gj}
  =
  \left\{
    \operatorname{SE}(\widehat{\beta}_{gj})^2+
    \widehat{\tau}_g^2
  \right\}^{-1},
  \qquad
  \widehat{\beta}_g
  =
  \frac{\sum_jw_{gj}\widehat{\beta}_{gj}}{\sum_jw_{gj}}.
  \label{eq:genere}
\end{equation}
BARCS-RE-EB (implementation label \texttt{BARCS-EB}) stabilizes the noisy
gene-specific heterogeneity estimate by shrinking it toward the pooled
screen-wide excess-$Q$ estimate $\tau_0^2$:
\begin{equation}
  \widetilde{\tau}_g^2
  =
  \frac{
    (m_g-1)\widehat{\tau}_g^2+d_0\tau_0^2
  }{
    (m_g-1)+d_0
  },
  \qquad d_0=4.
  \label{eq:geneeb}
\end{equation}
Both random-effects statistics are calibrated against prespecified
negative-control genes when at least ten are available and otherwise against a
robust whole-screen empirical null.  In contrast, BARCS-NORM obtains its
spread and reference degrees of freedom entirely within each gene.  These
three alternatives and BARCS-ST reuse exactly the same fitted guide
coefficients; only the guide-to-gene inference changes.

\paragraph{Guide-consistency statistic for the one-replicate boundary.}
When sample replication is insufficient, converting a guide statistic through
a low-degree-of-freedom $t$ distribution can erase useful ranking information.
We therefore prespecified a separate exploratory gene statistic, BARCS-GC,
for the $R=1$ CRISPulator analysis.  It uses the uncalibrated guide
coefficient and model-based standard error to estimate one shared gene effect.
With
\begin{equation}
  w_{gj}
  =
  \operatorname{SE}(\widehat{\beta}_{gj})^{-2},
\end{equation}
the inverse-variance estimate, its standard error, and raw Wald statistic are
\begin{equation}
  \widehat{\beta}_g
  =
  \frac{\sum_{j=1}^{m_g}w_{gj}\widehat{\beta}_{gj}}
       {\sum_{j=1}^{m_g}w_{gj}},
  \qquad
  \operatorname{SE}(\widehat{\beta}_g)
  =
  \left(\sum_{j=1}^{m_g}w_{gj}\right)^{-1/2},
  \qquad
  T_g
  =
  \frac{\widehat{\beta}_g}
       {\operatorname{SE}(\widehat{\beta}_g)}.
  \label{eq:guideconsistency}
\end{equation}
This is a direct fixed-effect estimate of the common guide coefficient, not a
Fisher or Stouffer combination of guide $p$-values.  Its additional assumption
is that the guide coefficients target one common gene effect; the reported
direction-agreement fraction exposes gross heterogeneity.

Because $T_g$ does not have a trustworthy sample-replicate reference
distribution at $R=1$, it is calibrated across genes.  Let $c_0$ be the
all-gene median and
\begin{equation}
  s_{\mathrm{MAD}}
  =
  \frac{\operatorname{median}_g|T_g-c_0|}
       {\Phi^{-1}(0.75)}.
  \label{eq:gcmad}
\end{equation}
When at least ten valid control genes are present, let $c_C$ be their median
and let
\begin{equation}
  s_C
  =
  \frac{Q_{0.95}\{|T_g-c_C|:g\in C\}}
       {\Phi^{-1}(0.975)}.
  \label{eq:gccontrol}
\end{equation}
The calibrated statistic and two-sided empirical-null $p$-value are
\begin{equation}
  Z_g^{\mathrm{GC}}
  =
  \frac{T_g-c_C}{\max(1,s_{\mathrm{MAD}},s_C)},
  \qquad
  p_g^{\mathrm{GC}}
  =
  2\Phi(-|Z_g^{\mathrm{GC}}|),
  \label{eq:gcz}
\end{equation}
followed by Benjamini--Hochberg adjustment.  If too few control genes are
available, $c_0$ replaces $c_C$ and $s_C$ is omitted.  At least three finite
guide coefficient estimates are required per gene.  The all-gene MAD assumes
that fewer than half of tested genes are active.

BARCS-GC measures reproducibility across independently designed
perturbations; it does not turn guides into independent biological samples.
Its $p$-value is consequently labeled empirical-null guide-consistency
evidence, not sample-level confirmation.  The ordinary BARCS result remains
the primary analysis when biological replicates are available, and
independent experimental replication remains necessary for confirmatory
claims.
